\usepackage[ngerman]{babel}
\usepackage[T1]{fontenc}
\usepackage{lmodern}
\usepackage[utf8]{inputenc}
\usepackage{hyperref}
\usepackage{amssymb}
\usepackage{amsmath}
% Dimensionen bitte nicht ändern.
\usepackage[left=2cm, right=2cm, top=2cm, bottom=2cm, bindingoffset=1cm, includeheadfoot]{geometry}
%Zeilenabstand bitte nicht ändern
\usepackage[onehalfspacing]{setspace}

\usepackage[backend=biber,style=numeric,]{biblatex}\addbibresource{literatur.bib}

% --- zusätzliche Pakete (Layoutmaße der Vorlage bleiben unverändert) --------
\usepackage{csquotes}
\usepackage{graphicx}
\usepackage{booktabs}
\usepackage{array}
\usepackage{xcolor}
\usepackage{float}
\usepackage{verbatim}
\usepackage{listings}
\usepackage{microtype}
\usepackage[font=small]{caption}

\hypersetup{colorlinks=true, linkcolor=black, citecolor=black, urlcolor=blue!60!black}

% --- Angaben zur Abgabe ----------------------------------------------------
\newcommand{\repourl}{https://github.com/yaniktfl/iv-repo}
\newcommand{\repobranch}{main}
\newcommand{\repocommit}{ce20021}

% --- Elm-Listings ----------------------------------------------------------
\lstdefinelanguage{Elm}{
  morekeywords={module,exposing,import,type,alias,let,in,if,then,else,case,of,port,as,where,True,False},
  morekeywords=[2]{Model,Msg,Maybe,Just,Nothing,List,Set,Dict,Float,Int,String,Bool,Html,Svg,Cmd,Sub,Result,Ok,Err},
  sensitive=true,
  morecomment=[l]{--},
  morecomment=[n]{\{-}{-\}},
  morestring=[b]",
}
\lstset{
  language=Elm,
  basicstyle=\ttfamily\footnotesize,
  keywordstyle=\color{blue!55!black}\bfseries,
  keywordstyle=[2]\color{teal!70!black},
  commentstyle=\color{gray}\itshape,
  stringstyle=\color{purple!60!black},
  showstringspaces=false,
  columns=fullflexible,
  breaklines=true,
  frame=single,
  rulecolor=\color{gray!50},
  framesep=4pt,
  xleftmargin=6pt,
  captionpos=b,
  aboveskip=0.8em,
  belowskip=0.5em,
  literate={ä}{{\"a}}1 {ö}{{\"o}}1 {ü}{{\"u}}1 {Ä}{{\"A}}1 {Ö}{{\"O}}1 {Ü}{{\"U}}1 {ß}{{\ss}}1
           {€}{{\texteuro}}1 {·}{{$\cdot$}}1 {→}{{$\rightarrow$}}1 {Ø}{{\O}}1 {–}{{--}}1,
}

% --- Abbildungsmakro -------------------------------------------------------
% #1 Dateiname ohne Endung in figures/, #2 Breite als Anteil von \textwidth,
% #3 Bildunterschrift, #4 Beschreibung für den Platzhalter.
\newcommand{\abb}[4]{%
  \begin{figure}[H]
    \centering
    \IfFileExists{figures/#1.png}
      {\includegraphics[width=#2\textwidth]{figures/#1}}
      {\fbox{\begin{minipage}[c][0.16\textheight][c]{0.92\textwidth}
        \centering\small
        \textbf{Platzhalter für \texttt{figures/#1.png}}\\[0.5em]
        #4
      \end{minipage}}}
    \caption{#3}\label{fig:#1}
  \end{figure}}
